%% ARXIV-ONLY fork of sections/overview.tex — arxiv.tex inputs this
%% file; main.tex keeps sections/overview.tex untouched. One question
%% down the use plane, then one growth step down the evolution plane —
%% the two planes narrated concretely before the theory begins.
%% Cherry-pick by diffing against sections/overview.tex.

\section{Overview: one question, one growth step}
\label{sec:overview}

Consider a small C program: a loop over an array with an index
computation, and a question --- can \texttt{x == 42} hold at the failure
label? \Cref{fig:run} draws the route this question takes; we narrate
it, then one growth step, and \Cref{sec:calculus} makes the
definitions precise.

\begin{figure}
\centering
\begin{tikzpicture}[x=1cm, y=1cm]
\node[lang] (c)    at (0.0, 0.0) {C};
\node[lang] (rv)   at (1.9, 0.0) {riscv};
\node[lang] (sail) at (3.0, 1.0) {sail};
\node[hub]  (bt)   at (4.1, 0.0) {btor2};
\node[hub]  (smt)  at (6.3, 0.0) {smtlib};
\draw[repr] (c) -- (rv);
\draw[chk] (rv) -- (bt);
\draw[chk] (rv) -- (sail);
\draw[chk] (sail) -- (bt);
\draw[pred] (bt) -- (smt);
\draw[->, >=stealth, dotted, thick]
  (smt.south) .. controls (4.6, -1.1) and (3.2, -1.1) .. (rv.south)
  node[midway, below, font=\scriptsize]
  {\textsc{sat} model $\to$ carry back $\to$ replay};
\end{tikzpicture}
\caption{One run down the spine. The question travels
C$\to$RISC-V$\to$BTOR2$\to$SMT-LIB (RISC-V reaching BTOR2 twice, the
manual-derived and Sail-derived branches whose agreement corroborates
both); a \textsc{sat} model is carried back hop by hop and
\emph{replayed} in the source-side interpreter --- the answer's
evidence, not the solver's say-so (\Cref{thm:existential}).}
\label{fig:run}
\end{figure}

\paragraph{The spine}
The program first moves from C to RISC-V. The translator on this edge is
a \emph{pinned} C compiler: a digest-locked gcc with recorded flags. We
can replay it bit-for-bit, but nobody can foresee or prove its output ---
its honest grade is $\trepr$, determinism without meaning. The RISC-V
program then moves to BTOR2, a bit-vector transition-system logic that
model checkers consume, and from there across a rule-for-rule bridge to
SMT-LIB, where bit-vector solvers decide the reachability query.

Each of these moves is a \emph{pair}: translator, source interpreter,
target interpreter, and a target-to-source interpreter that carries
behaviors backwards. When the RISC-V-to-BTOR2 translator runs, the
platform can immediately check it \emph{on this very program}: interpret
the RISC-V program directly (the left edge of a square), interpret the
BTOR2 translation and carry the resulting trace back to RISC-V
observables (the other three edges), and compare, step by step, on the
declared observables --- program counter, registers, halt flag. That
check passing is what the edge's $\tchk$ grade \emph{means}. Note
carefully what it covers: every hop \emph{from the compiled binary
onward}. The C hop itself has no square --- there is no C interpreter
and no ISA-to-C carry-back --- so the verified portion of the route
begins at the binary; the pinned compiler contributes reproducibility,
re-checked separately by a corpus-level differential against an
independent C verifier (\Cref{sec:instantiation}).

\paragraph{The branch}
RISC-V reaches BTOR2 twice. One translator was written against the prose
ISA specification; a second route factors through Sail, deriving its
semantics from the official RISC-V Sail model. On the segment where they
diverge --- from RISC-V to BTOR2 --- the two routes share no translator
and no carry-back; below the reconvergence they share the hub bridge,
so it is the diverse prefix that agreement corroborates
(\S\ref{sec:branching}). The
platform runs the question down both, and the answers must agree. If they
do, two independently derived encodings of RISC-V semantics corroborate
each other --- evidence strictly stronger than either route's own grade
(\Cref{lem:agree}). If they disagree, at least one route has a defect,
and the per-hop squares localize it: which hop, which step, which
observable (\Cref{lem:disagree,cor:localization}). The same dual-route
structure exists for AArch64 (Arm manual vs.\ Arm Sail model).

\paragraph{The answer comes home}
Suppose the SMT solver answers \textsc{sat}: the label is reachable, with
a model. The model is not the answer the user asked for --- it is a
valuation of SMT variables. The route's carry-backs now run in reverse:
the SMT model becomes a BTOR2 witness, the witness a RISC-V trace, the
trace an initial memory and register state --- concretely, the inputs of
the original C program. The platform then \emph{replays} those inputs in
the route's source-side interpreter --- for this spine, the shared
RISC-V interpreter on the compiled binary, the question's observable
read at the ISA level (\S\ref{sec:composition}) --- and watches the
failure state actually happen.
At that moment the user's trust obligations collapse: even if every
translator, both routes, and the solver were wrong, the replayed run
exhibits the behavior (\Cref{thm:existential}). Witness-carrying answers
are self-certifying; only the interpreter that replays them --- itself
differentially tested against an independently generated ISA simulator ---
remains to be believed.

Had the solver answered \textsc{unsat}, no replay could vouch for it.
That is where the grades, the branch agreement, the multi-engine
corroboration, and the independently checked certificates carry the
weight (\Cref{thm:universal}) --- and where the platform spends its
assurance budget.

\paragraph{Beyond one spine}
Nothing above is specific to C or RISC-V
(\Cref{fig:registry} draws the full registry graph). The same square contract admits
eBPF, Wasm, and EVM programs into the BTOR2 hub; Python functions and
chemical reaction networks reach the SMT-LIB hub directly; a SMILES
molecular graph reaches a molecular-formula target with no solver at all
(a \emph{compile pair} --- the calculus does not care that the languages
are far from computation, only that both have defined meaning). Each
pair declares what it preserves (its projection), what it rejects (typed,
enumerable $\unsupported$), and its grade; the registry composes routes
and reports each route's composed coverage, grade, and loss.

\paragraph{One growth step}
Now ask the platform something it cannot do: is this RISC-V program's
loop \emph{live} --- does it always eventually make progress? The
diagnosis (\texttt{gurdy why-not}) walks its five obstacles in order
and stops at the third: routes exist, the observables survive, but no
registered reasoning language expresses a liveness shape. The failure
is recorded in the platform's books --- the question verbatim, the
obstacle, the generation target it names (a liveness-to-safety
reduction on the bit-level hub, or a new temporal hub), an origin tag
--- and the recommendation board aggregates it beside every other
unmet demand. A generator drafts the brief citing those records; a
human reads the evidence and registers, or declines; a builder agent
implements against the contract; the gate admits or rejects; the
ratchet keeps every prior verdict standing; the question is re-asked.
That is one turn of the evolution plane, and no step of it answered a
question or was answered by one --- \emph{answers never write; growth
never answers}. The rest of
the paper makes each of these words precise, then measures the system.

\paragraph{The graph in the limit}
Nothing above says the graph is finished --- the architecture is built
so that it never has to be. When a question is unanswerable, exactly
one of four obstacles fails first, each named by one shipped call
(\texttt{gurdy why-not}): no route reaches a reasoning language; a
route exists but some hop's declared loss discards the observable the
question reads; no registered reasoning language expresses the
question's shape; or the solver exhausts its budget. The diagnosis
names the missing pair; the demand is recorded in the platform's
books and aggregated into recommendations (\S\ref{sec:books}); and
that makes a \emph{loop} well-defined: generate the recommended pair,
gate it, ratchet, re-ask. We leave the closed-loop composition as a
concept --- its stages are shipped, tested calls, and no autonomous
run is claimed or evaluated here --- but its limit
clarifies what the platform is \emph{for}. The answerable set is the
closure of a source language's questions under the registered
reductions; it only grows (\Cref{prop:ratchet}); and any question that
some finite chain of sound, deterministic reductions can bring to a
mechanized decision procedure is eventually answered, with quantified
trust. For a fixed finite-state program --- the compiled binary
above --- every question over kept observables is decidable in
principle, so growth for a \emph{given} program is a matter of
question shapes (new reasoning hubs) and of \emph{cost} (abstraction
pairs, \S\ref{sec:lax}), not connectivity. Trust grows on a separate
axis --- an added route answers no new question; it corroborates
(\S\ref{sec:branching}) --- and saturates at the supply of independent
semantic anchors, which no generator enlarges. Undecidability and
interpreter adequacy (\Cref{asm:adequacy}) stand at every graph size:
the limit is not an oracle but the reflection of everything mechanized
deciding knows how to answer, behind one interface, with every
answer's trust itself measured. The artifact's \texttt{POTENTIAL.md}
develops the limit and its boundaries in full.
